\begin{abstract}
Wildlife surveys ask \emph{how many animals are there}, but the computer-vision pipelines
built to answer that question are optimized as \emph{detectors}. We show that on thermal
road-transect video these are different problems, and we localize precisely where the
difference costs animals. Using a new corpus of 32 fully annotated FLIR thermal transects
(521,930 frames, 236 individually tracked white-tailed deer, 21,646 boxes), we benchmark 12
detectors and then decompose end-to-end counting error into three measurable stages: how
many animals are \emph{detected at all}, how many acquire their \emph{own} candidate track,
and how many are finally \emph{counted}. On held-out video the pipeline detects 98.8\% of
deer, gives 88.0\% a distinct track, and counts 62.7\%. Detection is not the bottleneck.
We then show something counter-intuitive and, we argue, generalizable: four independent
interventions that each \emph{improve} candidate generation --- lowering the detector
threshold, removing the tracker's initialization gate, loosening NMS, and adding appearance
re-identification --- each make the final count \emph{worse}, because the confirmation stage
cannot exploit a pool it did not shrink. Under three evaluation protocols, learned
confirmation heads (transformer, gradient boosting, logistic regression) all lose to a
three-parameter hand-tuned rule. We further show that the standard IoU$\geq$0.50 criterion
re-ranks detectors relative to a counting-appropriate criterion (8 of 12 models shift
$\geq$2 places), and that the detector fails at \emph{both} scale extremes --- the large end being a hole in the training distribution
rather than a capability limit. We release the dataset, annotations, and evaluation harness.
\end{abstract}
